\section{Does Agglomeration Cost the Instance Margin?}\label{sec:ablation}
Agglomerative backbones distil several heterogeneous teachers into one student
\cite{ranzinger2024amradio,heinrich2025radiov25}. The concern this raises for
re-identification is specific: classification asks whether a \emph{class} is separable, ReID
asks whether one \emph{individual} is separable from a visually near-identical other, and
averaging teacher objectives is a plausible way to lose exactly that margin while every
classification benchmark continues to look healthy.

The concern is testable here rather than by analogy, because C-RADIOv4's teachers are in the
encoder set. \Cref{tab:ablation} is the ablation: the student at two capacities, its literal
teacher SigLIP2-g at 1163M, a capacity-matched SigLIP2 at 428M, and the DINOv3
family bracketing the student from above and below. All rows use the ArcFace head, under which
every probe in the study converged, and all teacher rows use the \texttt{squash} geometry so
that \cref{sec:geometry}'s confound is not silently inside this table.

\begin{table*}[t]
\centering
\caption{The teacher ablation, ArcFace mAP. C-RADIOv4 is a student of the SigLIP2 and DINOv3
families; SigLIP2-g is the literal teacher. \emph{role} names that relation. Market-1501 is the
only column whose identities any head has seen. DINOv3 stands in for a teacher family, not for
the actual DINOv3-7B teacher, which does not fit on this study's GPU.}
\label{tab:ablation}
\small
% generated by tools/gen_tables.py from results/runs — edit the runs, not this
\begin{tabular}{lllrrrrrrr}
\toprule
encoder & M & role & Market & MSMT17 & CUHK03 & Occ-REID & CCVID & VRIC & VRAI \\
\midrule
C-RADIOv4-H & 653 & student & \textbf{0.709} & 0.165 & 0.153 & 0.646 & 0.609 & 0.191 & 0.228 \\
C-RADIOv4-SO400M & 431 & student & 0.696 & 0.162 & 0.140 & \textbf{0.651} & 0.625 & 0.197 & 0.207 \\
\midrule
SigLIP2-g & 1163 & teacher & 0.608 & \textbf{0.231} & \textbf{0.283} & 0.649 & 0.626 & 0.115 & 0.229 \\
SigLIP2-SO400M & 428 & teacher & 0.514 & 0.200 & 0.174 & 0.558 & 0.589 & 0.075 & 0.177 \\
\midrule
DINOv3-H+ & 840 & teacher family & 0.655 & 0.067 & 0.097 & 0.356 & 0.461 & \textbf{0.359} & \textbf{0.247} \\
DINOv3-L & 303 & teacher family & 0.607 & 0.069 & 0.075 & 0.332 & 0.434 & 0.282 & 0.160 \\
\midrule
CLIP-B/16 & 87 & control & 0.289 & 0.027 & 0.017 & 0.414 & \textbf{0.795} & 0.035 & 0.039 \\
\bottomrule
\end{tabular}

\end{table*}

\subsection{In-domain, the margin improves rather than erodes}
On Market-1501 the student beats its literal teacher, \studentMarket{} against
\teacherMarket{} mAP, despite the teacher having 1.8$\times$ its parameters and running
at a higher input resolution that \cref{sec:geometry} shows is an advantage rather than a
neutral difference.

More to the point, the advantage is \emph{larger} on the metric the concern is about. mINP
\cite{ye2022deep} measures the rank of the hardest true match, which is the instance margin
that distillation was predicted to destroy. In-domain the student's mINP is
\leadMarketMINP{} against the teacher's \teacherMarketMINP{}, a factor of \minpGap{}$\times$
where the mAP factor is smaller. Given labels in the target domain, agglomerative distillation
does not merely preserve the fine-grained margin --- it improves it.

\subsection{Off-domain, the advantage does not survive}
It is a narrower claim than it first appears, and the rest of \cref{tab:ablation} says why.
Market-1501 is the only dataset whose identities any head has seen. On the five transfer sets
the student never exceeds the better teacher family by more than the seed noise of
\cref{sec:setup}: it ties on Occluded-REID (\familyStudentOccluded{} against
\familySiglipOccluded{}) and CCVID (\familyStudentCcvid{} against \familySiglipCcvid{}), and
loses on CUHK03-NP, VRIC and VRAI. On CUHK03-NP it loses to a teacher 1.5$\times$
\emph{smaller} than itself.

So the property is not ``agglomeration wins on people'' --- CUHK03-NP is a person dataset and
falsifies that --- and it is not ``agglomeration wins'' either. It is that \textbf{the
student's gain is concentrated exactly where the probe had labels}, which is also why the
retention view of the same numbers (\cref{tab:retention}) is not independent evidence: high
in-domain performance with mediocre transfer \emph{is} low retention, and reporting both as
two findings would be counting one twice.

\begin{table*}[t]
\centering
\caption{Retention: target mAP divided by source mAP, for the one head fitted on
\texttt{market1501/train}. The absolute source column is printed beside the ratios on purpose
--- a retention ratio flatters a weak source model, and CLIP's Occluded-REID and CCVID entries
are exactly that trap.}
\label{tab:retention}
\scriptsize
\setlength{\tabcolsep}{3pt}
% generated by tools/gen_tables.py from results/runs — edit the runs, not this
\begin{tabular}{lrrrrrrr}
\toprule
encoder & Market (source) & MSMT17 & CUHK03 & Occ-REID & CCVID & VRIC & VRAI \\
\midrule
C-RADIOv4-H & 0.709 & 0.23 & 0.22 & 0.91 & 0.86 & 0.27 & 0.32 \\
C-RADIOv4-SO400M & 0.696 & 0.23 & 0.20 & 0.94 & 0.90 & 0.28 & 0.30 \\
SigLIP2-g & 0.608 & 0.38 & 0.46 & 1.07 & 1.03 & 0.19 & 0.38 \\
SigLIP2-SO400M & 0.514 & 0.39 & 0.34 & 1.09 & 1.15 & 0.15 & 0.35 \\
DINOv3-H+ & 0.655 & 0.10 & 0.15 & 0.54 & 0.70 & 0.55 & 0.38 \\
DINOv3-L & 0.607 & 0.11 & 0.12 & 0.55 & 0.71 & 0.46 & 0.26 \\
CLIP-B/16 & 0.289 & 0.09 & 0.06 & 1.43 & 2.75 & 0.12 & 0.14 \\
\bottomrule
\end{tabular}

\end{table*}

\subsection{The transfer loss has a shape, and it is predictable}
The interesting result is not that the student loses off-domain but \emph{where}. The two
teacher families disagree sharply and by domain: language-aligned features carry degraded
person crops, and self-supervised dense features carry vehicles. SigLIP2 leads DINOv3 by
\spreadCuhk{}$\times$ on CUHK03-NP and by \spreadMsmt{}$\times$ on MSMT17; DINOv3 leads SigLIP2
by \spreadVric{}$\times$ on VRIC. \Cref{tab:bracket} places the student inside that
disagreement.

\begin{table*}[t]
\centering
\caption{Best-of-family ArcFace mAP, with the teachers' spread and the student's position
inside it. \emph{spread} is the ratio of the better teacher family to the worse; \emph{vs.\
best} is the student's relative distance from the better one. The two datasets where the
teachers disagree by roughly 3$\times$ both cost the student roughly
45\%.}
\label{tab:bracket}
\scriptsize
\setlength{\tabcolsep}{3pt}
% generated by tools/gen_tables.py from results/runs — edit the runs, not this
\begin{tabular}{lrrrrcr}
\toprule
dataset & SigLIP2 & DINOv3 & C-RADIOv4 & spread & bracketed & vs.\ best \\
\midrule
Market & 0.608 & 0.655 & \textbf{0.709} & 1.08$\times$ & \loss{above} & +8.2\% \\
MSMT17 & 0.231 & 0.069 & \textbf{0.165} & 3.35$\times$ & yes & \loss{-28.7\%} \\
CUHK03 & 0.283 & 0.097 & \textbf{0.153} & 2.91$\times$ & yes & \loss{-45.9\%} \\
Occ-REID & 0.649 & 0.356 & \textbf{0.651} & 1.82$\times$ & \loss{above} & +0.2\% \\
CCVID & 0.626 & 0.461 & \textbf{0.625} & 1.36$\times$ & yes & \loss{-0.0\%} \\
VRIC & 0.115 & 0.359 & \textbf{0.197} & 3.12$\times$ & yes & \loss{-45.1\%} \\
VRAI & 0.229 & 0.247 & \textbf{0.228} & 1.08$\times$ & \loss{below} & \loss{-7.7\%} \\
\bottomrule
\end{tabular}

\end{table*}

The pattern is regular enough to state as a rule of thumb. Where the teacher families agree
--- Occluded-REID at \spreadOccluded{}$\times$, CCVID at \spreadCcvid{}$\times$, VRAI at
\spreadVrai{}$\times$ --- the student matches the better of them to within
8\% and usually to within noise. Where they disagree by roughly a factor of three
--- CUHK03-NP at \spreadCuhk{}$\times$ and VRIC at \spreadVric{}$\times$ --- the student falls
\deficitCuhk{}\% and \deficitVric{}\% short of the better teacher respectively. Two independent
datasets, one person and one vehicle, at nearly the same teacher spread, produce nearly the
same student deficit.

\textbf{This is a different failure than the one the literature warns about, and a more useful
one.} The cost of agglomeration here is not a destroyed instance margin --- \cref{sec:ablation}
opens by showing the margin improving where labels exist. It is \emph{regression toward the
teacher mean on any domain where one teacher is much better than the other}. That cost is
predictable in advance from a question a practitioner can answer without running the
experiment: do my images look more like what the language-aligned teacher was good at, or more
like what the self-supervised one was good at? Where the honest answer is ``strongly one of
them'', the blend is the wrong purchase and the better single teacher is available.

We report one prospective check on this. Both C-RADIOv4 rows for CUHK03-NP were forecast to
fall inside the teacher bracket before they were run, on the strength of the pattern in the
datasets already measured, and both did.

\subsection{A family signature: the angular margin does not transfer for DINOv3}
\Cref{tab:margin} showed that replacing the linear head with an ArcFace head helps almost
everywhere. The exceptions are concentrated in one family: in DINOv3's feature space, on
transfer sets, the angular margin loses. Both DINOv3 probes reach a perfect training accuracy
over the source identities, so this is not an optimisation failure --- the head fits, and what
it fits transfers worse than the plainer head does. The practical reading is uncomfortable and
worth stating plainly: \emph{fit an angular-margin head on your labelled domain} is good advice
for an agglomerative or language-aligned backbone and bad advice for a DINOv3 backbone deployed
cross-domain, which is the opposite of what a Market-1501-only study would conclude.

\subsection{One anomaly we cannot yet explain}
CCVID is the one column in \cref{tab:ablation} where the ordering does not resemble any other.
The 87M CLIP control, run under the \texttt{squash} geometry, scores \clipCcvid{} mAP
there, ahead of every larger encoder in the study including its own 1163M
language-aligned relative. Its retention on that dataset is correspondingly implausible
(\cref{tab:retention}). The result is reproducible from the cached features and the protocol
digest is shared with every other row in the column, so it is not an adapter fault of the kind
\cref{sec:substrate}'s guards are designed to catch; but we do not have a mechanism for it, the
tracklet pooling step is unique to this dataset, and we flag it as unresolved rather than
interpret it. It is listed among the open items in \cref{sec:limitations}.
